\section{Tokenized-to-Tokenized Transformations}
\label{sec:tok2tok}

This section covers two transformations that take place within the Tokenized carrier level: Narrative Abstraction (Narrative $\rightarrow$ Narrative) and Schematic Formalization (Narrative $\rightarrow$ Schematic). Both pathways share one source---the concrete trajectories produced during an agent's sequential decision-making (reasoning traces, tool-use trajectories, GUI interaction traces, or embodied execution records). The target artifact of Narrative Abstraction stays in natural-language form (reflections, rules, insights), which downstream agents read and reuse directly through language and multimodal understanding; the target artifact of Schematic Formalization carries an explicit structural constraint (code, workflows, graphs), whose structure downstream agents exploit mechanically through a parser, executor, or graph traversal. Both pathways have a clear boundary with Parametric Externalization: they take concrete agent interaction trajectories as input, with the model acting as abstractor, formalizer, or synthesizer; if an artifact is produced mainly by the model from a task description or internalized knowledge, with no concrete trajectory as input, it belongs to the latter.

\subsection{Narrative Abstraction: Narrative-to-Narrative Transformation}
\label{sec:narrative-abstraction}

Narrative Abstraction derives a new Narrative artifact from a raw Narrative trajectory through abstraction, compression, rewriting, or induction---the whole transformation completed within the Tokenized level. The transformed artifact takes forms such as reflection, summary, rule, insight, or skill; a later agent retrieves it and injects it into the decision context as in-context guidance, reusing it directly through language and multimodal understanding. We group methods by how the experience store is maintained: a Static Experience Store never revises an entry once it is written, whereas a Dynamic Experience Store involves rewriting, merging, or eviction.

\subsubsection{Static Experience Store}

Static Experience Store methods accumulate the products of experience abstraction as append-only entries. The concrete form of the store varies: some are temporary caches within a task, discarded when the task ends; others are persistent stores never revised after construction. What they share is that once an entry is written it is never edited or merged. We split them by the number of trajectories the abstraction draws on, into Single-Trajectory Experience Abstraction and Multi-Trajectory Experience Induction.

\paragraph{Single-Trajectory Experience Abstraction.} This class takes a single agent trajectory, or a local fragment of one, as input, and at the end of one execution or trial derives a Narrative artifact through reflection, summarization, critique, or failure explanation. Its premise is that a single trajectory already carries diagnostic decision information---a failed trajectory may expose mistaken assumptions, ineffective actions, or omitted constraints, and a successful trajectory may contain an effective decomposition strategy or environment-interaction pattern.

Reflexion~\cite{Shi23b} is an early representative of this paradigm. After each trial the agent generates a short reflection from the failed trajectory and the environment feedback, and prepends the error attribution and improvement suggestions to the next prompt. ReAP~\cite{Aza25} summarizes a single web-browsing trajectory into a Web-Reflection covering effective steps, site-functionality limits, shortcuts, and failure backtrack points. MemOrb~\cite{Hua25} compresses each dialogue turn and its tool-call trajectory, after the fact, into an Orb---a multi-field structure of observation, emotion, outcome, context, and strategic reflection---in a customer-service setting. Wu et al.~\cite{Wu26} deconstruct subquestion--subroutine pairs from large-scale reasoning trajectories synthesized by a teacher model, forming a static procedural datastore; at inference the model first states the current subquestion and then retrieves the matching subroutine as a hint. REFLECT~\cite{Liu23} hierarchically summarizes the multimodal observation stream of a single robot execution (RGB-D video, audio, and state records)---extracting a scene graph, event-level descriptions, and subgoal-level summaries in turn---and, when subgoal verification fails, generates a failure explanation that drives corrective replanning. WebCoach~\cite{Liu25e} compresses each completed web episode into a 3--5 sentence summary, retaining fields such as success state and failure mode, and stores it in an External Memory Store. In the code domain, SWE-Exp~\cite{Che25d} and MemGovern~\cite{Wan26ap} both distill structured experience entries from a repair loop: SWE-Exp extracts comprehension experience and modification experience separately from a single repair trajectory, forming experience-bank entries tagged with perspective, modification, and issue type; MemGovern distills an Experience Card from the closed-loop repair records of GitHub issues, PRs, and patches, with a checklist-style refine loop ensuring field quality.

\paragraph{Multi-Trajectory Experience Induction.} This class takes multiple agent trajectories as input and extracts a transferable Narrative artifact through comparison, clustering, or merging. Its premise is that the incidental factors of any single trajectory are identified and filtered out under cross-trajectory comparison, while recurring patterns or errors surface as more stable inductive results.

AutoGuide~\cite{Fu24} locates the first decision divergence point in pairs of good and bad trajectories for the same task, extracts a context-bound guideline around the divergent action, writes it in when-context then-action form, and indexes it by context. CFGM~\cite{Yan25} first derives coarse focus points from basic environment information and offline trial-and-error trajectories, then compares successful and failed trajectories to distill hybrid-grained tips. M2~\cite{Yan26b} distills topic-tagged insights from historical success logs into an external insight bank.

\subsubsection{Dynamic Experience Store}

Dynamic Experience Store methods treat the store as an object that keeps evolving with interaction. The system not only writes new experience but also rewrites, merges, and evicts existing entries, giving the store the ability to self-correct and de-duplicate over long-term use. We organize the discussion around four representative maintenance practices: Operation-Driven Editing, Survival-Driven Pruning, Skill Library Curation, and Multi-Tier Co-Evolution.

\paragraph{Operation-Driven Editing.} One line of work maintains a store of natural-language rules or insights and writes editing and voting as explicit memory operations. ExpeL~\cite{Zha23c} extracts transferable insights from a pool of successful trajectories and from failure--success pairs, and performs add, edit, upvote, and downvote on existing entries. H2R~\cite{Ye25b} splits hindsight reflection into high-level planning memory and low-level execution memory, supporting add, modify, upvote, and downvote. AutoManual~\cite{Che24} extracts rules differentially by three outcome types---direct success, indirect success, and failure---and, through dedicated consolidation and formulation stages, merges redundancy into a Markdown manual. ACE~\cite{Zha25f} has a Reflector extract insights, failure modes, and strategy bullets from the reasoning trajectory, and a Curator merge them as deltas into a Context Playbook, with semantic deduplication and pruning of low-value entries. Dynamic Cheatsheet~\cite{Suz25} has a Curator assess solving trajectories for usefulness and generalizability, distill the transferable content into a cheatsheet, and rewrite, delete, merge, and compress entries to control length.

\paragraph{Survival-Driven Pruning.} Another line drives the store toward convergence through replacement, eviction, or survival scoring. Darwinian Memory~\cite{Mi26b} writes GUI experience as memory units indexed by precondition and goal, and regulates the store through evolutionary replacement, survival-based pruning, and reputation-style inhibition. R2D2~\cite{Hua25e} truncates corrections out of historical web trajectories and generates replayable corrective reflections, continually replacing weaker old experience through a replay buffer and reflective update. Remember Me, Refine Me~\cite{Cao25} extracts keypoint-level procedural memory from a single trajectory, tracks retrieval and success utility over time, and performs verification, deduplication, and deletion of low-value memory. Mistake Notebook Learning~\cite{Su25} takes a batch of failure samples as the update unit, clusters them by subject and then generates mistake notes, accepting an update only after a post-update evaluation validates the batch as a whole.

\paragraph{Skill Library Curation.} Work centered on skills absorbs each new experience into a continually evolving skill library. Trace2Skill~\cite{Ni26} has parallel analysts propose local patches for individual trajectories, then synthesizes transferable skills through hierarchical merge and conflict-free consolidation, with existing skills further rewritten in a deepening mode. AutoSkill~\cite{Yan26} extracts skill candidates from user interaction signals, with a judge deciding add, merge, or discard. SkillClaw~\cite{Ma26b} aggregates real usage sessions across many users and agents, has an Agentic Evolver decide create, refine, or skip, and synchronizes verified skills back to a shared SkillHub.

\paragraph{Multi-Tier Co-Evolution.} A further line maintains several evolving stores with different semantic roles in parallel. Nemori~\cite{Ma25} uses the prediction error between an episode and an anticipatory schema to judge the retention value of information, distilling it into semantic knowledge, and fuses consecutive episodes into a narrative episodic memory; both tiers perform merge and conflict replacement. CLIN~\cite{Maj23} extracts causal language statements from each trial and summarizes a cross-episode meta-memory from high-quality memory, generating an updated version after each trial. CER~\cite{Liu25} extracts both local environment dynamics and higher-level decision-making skills from each trajectory, referencing existing content explicitly on write to avoid duplication.

\begin{table*}[t]
\centering
\caption{Overview of Narrative Abstraction methods.}
\label{tab:narrative-abstraction}
% Column definitions:
%   Store Type -- maintenance mode of the store (whether entries are revised after being written): {Static, Dynamic}.
%   Abstraction Source -- number of trajectories a single abstraction draws on: {Single-trajectory, Multi-trajectory}.
%   Update Mechanism -- how entries change over time: {append-only, operation-edit, survival-prune, skill-curate, multi-tier}.
%   Artifact Form -- concrete form of the product: {reflection, summary, rule, insight, skill, multi-field record, procedural pair}.
%   Retrieval Mechanism -- how entries are selected into the decision context: {always-on prepend, embedding similarity, keyed lookup, graph traversal, hierarchical drilldown, task-routed}.
%   Domain -- experimental setting or target domain: {text, web, GUI, mobile, embodied, code, multi-agent, customer-service}.
\scriptsize
\begin{tabular}{@{}lllllll@{}}
\toprule
Work & Store Type & Abstraction Source & Update Mechanism & Artifact Form & Retrieval Mechanism & Domain \\
\midrule
% Entries to be added.
\bottomrule
\end{tabular}
\end{table*}

\subsubsection{Discussion}

A static store is simple to implement, predictable in behavior, and cheap at run time, but experience quality is locked at write time---a mis-written reflection keeps affecting later retrieval. A dynamic store converges over long-term use through merge, replacement, and pruning, and is better at filtering the incidental noise of a single abstraction, at the cost of extra machinery for conflict resolution, staleness judgment, and version management; its behavior is harder to reproduce, and the eviction policy itself may delete entries that are still useful. Within the static class, single-trajectory abstraction has low latency, small data requirements, and good cold-start behavior, which suits narrow task distributions that need fast adaptation; multi-trajectory induction gains stronger generalization and noise resistance from cross-trajectory comparison, at the cost of weak cold-start behavior, the possibility that different trajectories support conflicting conclusions, and a risk of over-induction---wrongly elevating a local regularity of a specific environment into a general rule.

\subsection{Schematic Formalization: Narrative-to-Schematic Transformation}
\label{sec:schematic-formalization}

Schematic Formalization derives, from a raw Narrative trajectory, a Schematic artifact with explicit structure that can be parsed, executed, or traversed. We group methods by the main semantic function the Schematic artifact serves in later agent decisions, into three classes: Programmatic Skill Construction, Procedural Workflow Induction, and Structured Memory Graph Construction.

\subsubsection{Programmatic Skill Construction}

Programmatic Skill Construction builds callable code skills from concrete agent trajectories, so that an agent invokes them directly as high-level action units in later tasks rather than replanning from low-level actions. The source experience is usually a successful or filtered execution trajectory, and the product can be executed directly by an executor, runtime, or browser-automation framework.

Voyager~\cite{Wan23c} repeatedly generates, executes, and revises action programs in Minecraft through iterative prompting; only a program that passes environment feedback, execution checks, and self-verification is written to the skill library as a Mineflayer JavaScript program. ASI~\cite{Wan25d} starts from episodes that an evaluator judges successful during the agent's online solving, cleans out the erroneous steps, has an induction module induce typed Python functions, and rewrites the original trajectory into a version that calls the new skill before re-executing---only a skill that still reproduces success after the rewrite is admitted. SkillWeaver~\cite{Zhe25c} lets a web agent practice autonomously under an automatic curriculum, has a reward model identify successful behavior, hands the state--action sequence to an LLM to generate a Playwright API with typed parameters, and then polishes it continually with unit tests, static analysis, and environment feedback. WebXSkill~\cite{Wan26d} extracts action subsequences from trajectories synthesized by a stronger teacher model on WebArena and WebVoyager, abstracts them into JSON skills with typed parameters and step-level guidance, and organizes them into a skill graph keyed by URL pattern.

\subsubsection{Procedural Workflow Induction}

Procedural Workflow Induction induces reusable task-process structures from multi-step agent trajectories; the artifact takes the form of a workflow, hierarchical procedure, or task tree that encodes phase division, step order, dependencies, and pre/post conditions. The premise is that the reusable experience of many complex tasks lies in how the execution process is organized---successful trajectories often contain stable phase patterns, and failed trajectories may expose a missing step order or dependency check.

AWM~\cite{Wan24} extracts recurring sub-routines from the typical experience of a web agent; the product is a workflow trajectory composed of environment descriptions, reasoning, and action programs. WorkflowGen~\cite{Wei26} extracts experience at both node and workflow granularity, clustering structurally identical trajectories into a generalized workflow template. HMT~\cite{Tan26} and MACLA~\cite{For25} both segment trajectories and then perform hierarchical abstraction with explicit condition annotation: HMT builds an intent/stage/action three-level tree with pre/postconditions at the stage level; MACLA segments a trajectory into atomic procedures (the four-tuple of goal, precondition, action sequence, and postcondition) and abstracts the procedure sequence into a meta-procedure with continue/skip/repeat/abort control logic. MOBIMEM~\cite{Liu25i} distills mobile execution experience and user corrections into multi-level experience templates with parameter slots, and revises stale experience continually through exception-driven updates and rollback-based expiry detection. A2Flow~\cite{Zha25c} starts from expert demonstrations and task-resolution data, turns case-level demonstrations into code-form workflows and initial operators, and then uses MCTS to optimize the overall process over a search space of nodes, edges, and operators. FlowMind~\cite{Liu26} adopts a two-stage Execute--Summarize scheme: the model first completes the task with tools and records the execution trajectory, then, in a separate summarization stage, extracts the key steps, dependencies, and control flow into a workflow graph.

\subsubsection{Structured Memory Graph Construction}

Structured Memory Graph Construction turns the entities and relations in an agent trajectory into traversable, updatable graph-structured memory. Much agent experience fits a relational structure better than linear text---environment objects have spatial relations, actions change state, historical events have temporal or causal links, and multi-agent interactions involve roles, information flow, and dependencies. Stored only as raw logs or natural-language notes, these relations leave a later agent unable to do precise retrieval, local update, and multi-hop reasoning.

AriGraph~\cite{Ano24} and A-MEM~\cite{Xu25b} both turn a single-step interaction into a semantically structured atomic unit and establish relations between old and new units: AriGraph extracts a semantic triplet from each textual observation into semantic memory and deletes overturned old relations when the environment state changes; A-MEM turns each interaction into an atomic note with keywords, tags, and an embedding, establishes semantic relations between old and new notes through link generation, and lets new experience trigger rewrites of old notes. G-Memory~\cite{Zha25} targets multi-agent systems, organizing the atomic statements of a collaboration into an interaction graph, instantiating each whole query as a query node, and having an LLM summarize insight nodes, forming a three-tier interaction/query/insight graph. MobileGPT~\cite{Lee23} transcribes completed mobile-app tasks into a task/sub-task/primitive-action hierarchy through an Explore--Select--Derive process, and organizes page--subtask relations into a transition graph. Environment Maps~\cite{Fen26} compiles browsing trajectories, the DOM, and screenshots into a JSON environment map containing context nodes, parameterized action nodes, workflow order, and tacit knowledge. BrainMem~\cite{Ma26} maintains a dual graph in embodied settings---a Trajectory KG (action--state transitions) and a Spatial KG (room--object spatial relations); after an episode, an Experience Agent distills successful trajectories into general patterns and compresses failed trajectories, via guided retry, into symbolic criteria.

\begin{table*}[t]
\centering
\caption{Overview of Schematic Formalization methods.}
\label{tab:schematic-formalization}
% Column definitions:
%   Artifact Category -- semantic-function class in this subsection: {Programmatic Skill, Procedural Workflow, Structured Memory Graph}.
%   Artifact Form -- concrete structural form: {code skill, JSON skill, workflow trajectory, workflow template, hierarchical tree, transition graph, knowledge graph, environment map}.
%   Downstream Mechanism -- how the downstream agent uses the artifact: {execute, parse-and-instantiate, traverse}.
%   Retrieval Mechanism -- how the target entry is selected: {embedding similarity, keyed lookup, URL pattern matching, graph traversal, task-routed, hierarchical drilldown}.
%   Hierarchical Structure -- whether the artifact has multi-level abstraction: {flat, multi-level}.
%   Domain -- experimental setting or target domain: {Minecraft, web, mobile, multi-agent, embodied, code}.
\scriptsize
\begin{tabular}{@{}lllllll@{}}
\toprule
Work & Artifact Category & Artifact Form & Downstream Mechanism & Retrieval Mechanism & Hierarchical Structure & Domain \\
\midrule
% Entries to be added.
\bottomrule
\end{tabular}
\end{table*}

\subsubsection{Discussion}

Programmatic Skill Construction has the advantages of executability, verifiability, and composability: an artifact's validity can be judged from its run result, and several skills can be composed within a complex task. The cost is a high sensitivity to the stability of the environment interface, and a skill library that keeps growing needs accompanying selection, deduplication, versioning, and stale detection.

Procedural Workflow Induction has the advantages of interpretability and controllability: a workflow artifact is shorter and more abstract, expresses step dependencies and execution order explicitly, and makes the key phases of a task easy to locate. The main challenge is controlling the abstraction granularity---too specific and it does not generalize, too abstract and it loses key context---together with workflow drift: once the environment state changes, the original workflow may no longer hold, which calls for consistency checking, precondition verification, and an update mechanism.

A Structured Memory Graph has the advantage of supporting relation-level memory reuse; its graph structure makes local update convenient and suits the long-horizon planning and state tracking of embodied agents, web agents, and multi-agent systems. The difficulty is that graph extraction easily produces wrong entities and wrong relations, and a graph memory that grows over time develops memory bloat, stale information, and retrieval noise, which usually calls for accompanying graph update, conflict resolution, pruning, and confidence estimation.

The three artifact types are functionally complementary. Relative to a Programmatic Skill, a workflow sits at a coarser granularity---one skill can be a workflow node, and one workflow can orchestrate several skills, tools, or subtasks. Relative to a Procedural Workflow, a Structured Memory Graph encodes \emph{what exists in the environment, how state changes, and how historical experiences relate} rather than \emph{how a task should be executed}: the former uses nodes and edges for entities, states, events, and spatial or interaction relations, the latter for step order, control flow, or task dependencies.
